\begin{titlepage}
\vspace*{28mm}
{\Huge\bfseries\color{accent} stmlab + stmgui\par}
\vspace{4mm}
{\Large The complete Igor STM setup in Python, with its panels\par}
\vspace{14mm}
\rule{\linewidth}{0.6pt}
\vspace{4mm}

{\large Igor's three panels, nine graph windows and two macros, rebuilt on
top of the eight translated experiments --- and the reasoning behind every
control, every file, and every deviation from the original.\par}

\vspace{4mm}
\rule{\linewidth}{0.6pt}

\vspace{18mm}
\begin{keybox}
\textbf{Read this first.} This manual has three parts.

\textbf{Chapter 05, the structure map}, comes first and stands alone: the
whole code on a few pages --- what every file is, which file calls which
function in which other file (extracted from the source, not written), one
button click followed top to bottom, and one break-junction trace. Read it
before anything else; return to it whenever a chapter loses you.

\textbf{Part I, the GUI} (chapters 40--51), is new to this folder. Chapter 40
is the tour; chapter 41 is the lab flow you know from Igor, button by
button --- Start Writing, Approach, Find Offset, Find Suppress, Start
Measurement --- and is the one to read before the first run. Chapters
42--45 are the panels and the graphs, control by control. Chapter 47 is
what gets written to disk. Chapter 48 is how the code is built, for the day
something needs changing. Chapters 50--51 are the quick reference and the
complete binding table: every Igor global and where it lives now.

\textbf{Part II, the core} (chapters 00--30), is the stmlab manual: the
platform, one chapter per experiment, the support systems and the
command-line quick reference. It is unchanged from
\texttt{STMLAB\_Manual.pdf}; the GUI calls exactly this code. Chapter
numbers in print are the file numbers in \texttt{manual/}, which is why
they are not consecutive.

\medskip
Everything runs against a simulated junction with \texttt{-{}-simulate}.
Every figure in Part I was produced by the GUI itself, on the simulator.
\end{keybox}

\vfill
\begin{warnbox}\small
\textbf{Trust order.} The constant-bias path carries the reference tests and
the 1~\G{} self-check on the simulator. The four ramp modes, both campaigns
and the support modules are tested end to end on the simulator, through the
GUI as well as the command line. Every \texttt{nidaqmx} call, the Keithley
over GPIB, the second card and the electrochemical cell have \textbf{never
met their instruments}. The GUI changes none of this: it calls the same
functions, and is exactly as trustworthy as they are.

Simulate first. Scope second. Tip last.
\end{warnbox}

\vspace{4mm}
{\small \texttt{python\_code\_entire\_igor\_w\_GUI/}\quad$\cdot$\quad
3 panels, 9 graphs, 8 experiments\quad$\cdot$\quad 121 tests\quad$\cdot$\quad
translated from \texttt{igor\_code/}}
\end{titlepage}

\tableofcontents
